\scanpage{10}
9°) \uline{Th. de spécialisation} Voici \unsure{tjrs} une forme précise : $f:X\longrightarrow Y$ morphisme \uline{propre} et \uline{lisse}, $F$ faisceau de torsion loc. libre sur $X$, premier aux car. résiduelles de $X$, alors les $R^q f_*(F)$ [dont les fibres sont les coh. des fibres \ldots] sont loc. libres.

10°) \uline{Dim.} \uline{cohomologique}.

$X$ de type fini sur $k$ sép. clos, $n=\dim X$, alors
\[
H^i(X,F)=0\quad\text{pour }i\unsure{>|\geq}2n,
\]
($F$ faisceau de torsion sur $X$).

Si $X$ est \uline{affine}, on a $\unsure{=|\simeq}\ H^i(X,F)=0$ pour $i>n$ (th. du type Lefschetz).

La dimension ``topologique'' de $X$ \unsure{passe} donc pour être en général le double de la dimension de Krull --- il y a un abaissement dans le cas affine, qui peut donner des idées fausses.
\mspage{9}
